\documentclass[12pt,a4paper]{article}

% --- بسته‌های مورد نیاز ---
\usepackage{amsmath}
\usepackage{amssymb}
\usepackage{booktabs}
\usepackage{graphicx} % برای افزودن تصاویر
\usepackage{float}    % برای ثابت کردن جایگاه تصاویر در متن با دستور [H]

% --- تنظیمات زبان و فونت ---
\usepackage{xepersian}
\settextfont{B-NAZANIN.TTF} % فونت اصلی فارسی
\setlatintextfont[Scale=0.85]{Times New Roman}

% --- تنظیمات بولت‌پوینت‌ها برای جلوگیری از نمایش مربع ---
\renewcommand{\labelitemi}{$\bullet$}
\renewcommand{\labelitemii}{$\circ$}

\title{معرفی شبکه‌های عصبی پیش‌خور (\lr{FFN})}
\author{}
\date{}

\begin{document}
\maketitle

\section{شبکه \lr{FFN} چیست؟}

شبکه عصبی پیش‌خور یا \lr{FFN} (مخفف \lr{Feed Forward Neural Network})، ساده‌ترین و پایه‌ای‌ترین نوع شبکه‌های عصبی مصنوعی (\lr{ANN}) است. دلیل نام‌گذاری آن به «پیش‌خور» این است که اطلاعات در آن فقط در یک جهت (رو به جلو) حرکت می‌کنند و هیچ‌گونه حلقه (\lr{Loop}) یا بازگشتی در شبکه وجود ندارد.

یک شبکه \lr{FFN} از سه بخش اصلی تشکیل شده است:

\begin{itemize}
    \item \textbf{لایه ورودی (\lr{Input Layer}):} داده‌های خام را دریافت می‌کند (تعداد نورون‌ها برابر با تعداد ویژگی‌های داده است).
    \item \textbf{لایه‌های پنهان (\lr{Hidden Layers}):} پردازش اصلی و استخراج الگوهای غیرخطی در این لایه‌ها انجام می‌شود. یک شبکه می‌تواند یک یا چندین لایه پنهان داشته باشد. به شبکه‌هایی که بیش از یک لایه پنهان دارند، شبکه‌های عصبی عمیق (\lr{Deep Neural Networks}) می‌گویند.
    \item \textbf{لایه خروجی (\lr{Output Layer}):} نتیجه نهایی (مثل پیش‌بینی یک عدد یا دسته‌بندی یک تصویر) را تولید می‌کند.
\end{itemize}

% --- نحوه قرار دادن شکل (اختیاری برای درک بهتر لایه‌ها) ---
\begin{figure}[H]
    \centering
    % نام فایل تصویر ساختار شبکه عصبی خود را به جای ffn_diagram.png قرار دهید
    % \includegraphics[width=0.7\textwidth]{ffn_diagram.png} 
    \caption{معماری پایه یک شبکه عصبی پیش‌خور (\lr{FFN}) شامل لایه‌های ورودی، پنهان و خروجی.}
    \label{fig:ffn}
\end{figure}

\end{document}